% Part: turing-machines
% Chapter: machines-computations
% Section: church-turing-thesis

\documentclass[../../../include/open-logic-section]{subfiles}

\begin{document}

\olfileid{tur}{mac}{ctt}
\olsection{أطروحة تشرتش--تورنغ}

يفترض في آلات تورنغ أن تكون بديلًا دقيقًا لمفهوم الإجراء الفعال. وقد
رأى تورنغ أن كل من يفهم مفهوم الإجراء الفعال ومفهوم آلة تورنغ سيجد
حدسيًا أن كل ما يمكن فعله بإجراء فعال يمكن أن تفعله آلة تورنغ. ويعزز
هذه الدعوى أن البدائل الدقيقة الأخرى كلها التي اقترحت لمفهوم الإجراء
الفعال اتضح أنها مكافئة امتداديًا لمفهوم آلة تورنغ؛ أي إنها تستطيع
حساب مجموعة الدوال نفسها بالضبط. وتسمى هذه الدعوى \emph{أطروحة
تشرتش--تورنغ}.

\begin{defn}[أطروحة تشرتش--تورنغ]
تنص \emph{أطروحة تشرتش--تورنغ} على أن كل ما يمكن حسابه بإجراء فعال
قابل للحساب بتورنغ.
\end{defn}

يحتج بأطروحة تشرتش--تورنغ بطريقتين. والنوع الأول من استعمالها ذريعة
للكسل: فلنفترض أن لدينا وصفًا لإجراء فعال يحسب شيئًا ما، وليكن مكتوبًا
بـ«شفرة كاذبة». عندئذ يمكننا الاستناد إلى أطروحة تشرتش--تورنغ لتسويغ
الدعوى بأن آلة تورنغ ما تحسب الدالة نفسها، وإن لم نكن قد أنشأناها
بالفعل.

أما الاستعمال الآخر لأطروحة تشرتش--تورنغ فأهم فلسفيًا. إذ يمكن إثبات
وجود دوال لا تستطيع آلات تورنغ حسابها. ومن ذلك، مع استعمال أطروحة
تشرتش--تورنغ، يمكن استنتاج أنه لا يمكن حسابها حسابًا فعالًا بأي إجراء
كان. فلو وجد إجراء كهذا، للزم من أطروحة تشرتش--تورنغ وجود آلة تورنغ
له. ومن ثم، إذا أثبتنا عدم وجود آلة تورنغ تحسبها، لم يمكن أن يوجد
إجراء فعال أيضًا. ويحتج بأطروحة تشرتش--تورنغ، على وجه الخصوص، للقول
إن ما يسمى مسألة التوقف لا تعجز آلات تورنغ عن حلها فحسب، بل يستحيل
حلها حلًا فعالًا على الإطلاق.

\end{document}
